\section*{Заключение}
\addcontentsline{toc}{section}{Заключение}

В ходе выполнения летней практики была подготовлена учебная среда для моделирования работы очереди вычислительного кластера. На локальной машине были созданы пять виртуальных машин с openSUSE Leap: управляющий узел \texttt{mgmt} и вычислительные узлы \texttt{cn01}--\texttt{cn04}. Для узлов настроены статические IP-адреса, разрешение имён, SSH-доступ по ключам и общая конфигурация MUNGE.

На кластере установлен SLURM. Конфигурационный файл сформирован через официальный конфигуратор и синхронизирован между всеми узлами. Управляющий узел запускает службу \texttt{slurmctld}, вычислительные узлы запускают службу \texttt{slurmd}. Работоспособность кластера проверялась командами \texttt{sinfo}, \texttt{squeue} и запуском пакетных заданий.

Для индивидуального варианта был подготовлен фрагмент исходного набора данных с UID 50728 и именем задания \texttt{qe7}. По выбранным данным рассчитаны параметры распределения фактического времени выполнения и ошибки прогноза времени. Генератор формирует пакетные сценарии SLURM, выбирает количество узлов, рассчитывает плановое и фактическое время задания и отправляет задания в очередь. Аналитический модуль строит графики исходного распределения, сгенерированного распределения и фактических результатов выполнения.

Последний сформированный манифест содержит 500 заданий. Распределение числа узлов по заданиям составило 118 заданий на один узел, 245 заданий на два узла, 124 задания на три узла и 13 заданий на четыре узла. При выбранном масштабе времени медианное плановое время составляет 972 секунды, медианное время \texttt{sleep} составляет 74 секунды. Оценка длительности полного запуска на четырёх вычислительных узлах составляет около 6,5 часов, а ожидаемое отклонение фактического времени от планируемого при накладных расходах порядка 1,5 секунды на задание составляет около 2\%.

Таким образом, поставленная задача выполнена: учебный кластер SLURM развёрнут, генератор пользовательских заявок реализован, исходные данные обработаны, а результаты генерации и выполнения представлены в виде численных характеристик и графиков.
